\subsubsection{Ablation Study}

Để đánh giá một cách toàn diện mức độ đóng góp của từng nhóm đặc trưng
và từng họ hàm cơ sở, chúng tôi thực hiện ablation study bằng cách lần lượt
loại bỏ từng thành phần trong mô hình và quan sát sự thay đổi của sai số
bình phương trung bình (MSE) trên tập validation.

Mục tiêu chính của thí nghiệm này gồm:
(i) xác định các đặc trưng quan trọng nhất,
(ii) đánh giá vai trò của từng loại basis function,
(iii) phân tích mức độ nhạy của mô hình dưới các perturbation khác nhau.

%------------------------------------------------
\paragraph{1. Biểu diễn trực quan kết quả Ablation}

\begin{figure}[htbp]
    \centering
    \includegraphics[width=0.95\linewidth]{ablation_study.png}
    \caption{Ablation Study: Feature \& Basis Importance (MSE trên validation, thang log)}
    \label{fig:ablation_plot}
\end{figure}

Hình \ref{fig:ablation_plot} thể hiện kết quả ablation study dưới dạng
biểu đồ cột, trong đó:

\begin{itemize}
    \item Màu xanh lá: các mô hình đơn lẻ (individual models)
    \item Màu xanh dương: các mô hình loại bỏ từng thành phần (ablation models)
    \item Màu đỏ: mô hình đầy đủ (full model)
\end{itemize}

Do độ chênh lệch giữa các giá trị MSE rất lớn (từ $10^{-1}$ đến $10^{4}$),
trục tung được biểu diễn theo thang logarit ($\log_{10}$) nhằm đảm bảo khả năng
so sánh trực quan giữa các cấu hình.

%------------------------------------------------
\paragraph{2. Bảng tổng hợp kết quả Ablation Study}

Bảng \ref{tab:ablation_summary} tổng hợp toàn bộ giá trị MSE tương ứng với từng cấu hình mô hình.

\begin{table}[htbp]
\centering
\renewcommand{\arraystretch}{1.2}
\begin{tabular}{|l|c|}
\hline
\textbf{Model / Configuration} & \textbf{Validation MSE} \\
\hline
RBF               & 0.6379 \\
Fourier           & 0.4199 \\
Polynomial        & 65.1358 \\
NO\_RBF           & 0.4939 \\
NO\_Fourier       & 1767.3249 \\
NO\_Polynomial    & 0.4175 \\
NO\_Economic      & 255.3254 \\
NO\_Density       & 0.4427 \\
NO\_Spatial       & 605.1461 \\
NO\_MedInc        & 807.2610 \\
NO\_HouseAge      & 8.0081 \\
NO\_AveRooms      & 189.1508 \\
NO\_AveBedrms     & 90.5460 \\
NO\_Population    & 60.2631 \\
NO\_AveOccup      & 0.4718 \\
NO\_Latitude      & 187.5998 \\
NO\_Longitude     & 161.8832 \\
FULL              & 78.5531 \\
\hline
\end{tabular}
\caption{Tổng hợp kết quả Ablation Study trên tập validation (MSE)}
\label{tab:ablation_summary}
\end{table}

%------------------------------------------------
\paragraph{2. Phân tích tổng quan hành vi mô hình}

Quan sát từ biểu đồ, có thể rút ra một số đặc điểm quan trọng:

\begin{itemize}
    \item Nhóm mô hình Fourier và NO\_POLY đạt MSE thấp nhất ($\approx 0.41$--$0.42$),
    cho thấy đây là vùng tối ưu về mặt biểu diễn.

    \item Polynomial basis thể hiện hành vi cực kỳ không ổn định với MSE tăng mạnh,
    phản ánh sự bùng nổ sai số do tính chất ill-conditioned của không gian đa thức.

    \item Các cấu hình loại bỏ Fourier (NO\_Fourier) dẫn đến MSE tăng đột biến
    ($> 10^3$), cho thấy Fourier basis đóng vai trò cốt lõi trong việc tái tạo
    cấu trúc dữ liệu.

\end{itemize}

%------------------------------------------------
\paragraph{3. Phân tích theo nhóm đặc trưng (Feature Importance)}

Từ các cấu hình loại bỏ từng feature, ta có thể xếp hạng mức độ quan trọng
của các đặc trưng như sau:

\[
\begin{aligned}
\text{MedInc} \gg \text{Spatial Features (Latitude, Longitude)} \gg \\
\text{HouseAge} > \text{Population} > \text{Rooms Features} > \text{AveOccup}
\end{aligned}
\]

Cụ thể:

\begin{itemize}
    \item \textbf{MedInc (Median Income)}:
    khi loại bỏ làm MSE tăng lên $807.26$, chứng tỏ đây là biến có sức dự báo mạnh nhất,
    phản ánh mối quan hệ kinh tế trực tiếp đến giá nhà.

    \item \textbf{Spatial features (Latitude, Longitude)}:
    khi loại bỏ gây tăng MSE lớn ($\approx 160$--$187$),
    cho thấy giá nhà phụ thuộc mạnh vào vị trí địa lý (spatial dependency).

    \item \textbf{Population, AveRooms, AveBedrms}:
    ảnh hưởng trung bình, phản ánh cấu trúc khu dân cư.

    \item \textbf{AveOccup}:
    gần như không ảnh hưởng ($0.4718$), cho thấy biến này ít thông tin dự báo.
\end{itemize}

%------------------------------------------------
\paragraph{4. Phân tích theo basis function (Model Sensitivity)}

Xét theo nhóm hàm cơ sở:

\begin{itemize}
    \item \textbf{Fourier Basis}:
    ổn định nhất, đạt MSE thấp nhất ($0.4199$),
    chứng tỏ dữ liệu có cấu trúc mang tính chu kỳ/tần số thấp.

    \item \textbf{RBF Basis}:
    hoạt động tốt ($0.6379$), nhưng phụ thuộc vào phân bố center,
    do tính chất local approximation.

    \item \textbf{Polynomial Basis}:
    mất ổn định nghiêm trọng ($65.13$),
    phản ánh hiện tượng numerical explosion khi tăng bậc đa thức.
\end{itemize}

Đặc biệt, cấu hình NO\_Fourier làm MSE tăng lên hơn $10^3$,
chứng minh Fourier basis là thành phần không thể thiếu trong việc
mô hình hóa cấu trúc dữ liệu.

%------------------------------------------------
\paragraph{5. Hiện tượng đáng chú ý: FULL model kém hiệu năng}

Một kết quả đáng chú ý là mô hình FULL có MSE cao bất thường ($78.55$),
trong khi nhiều cấu hình rút gọn lại đạt MSE thấp hơn rất nhiều.

Điều này có thể được giải thích bởi:

\begin{itemize}
    \item \textbf{Over-parameterization}: số lượng feature quá lớn gây nhiễu gradient
    \item \textbf{Feature interaction không kiểm soát}: các basis kết hợp gây collinearity
    \item \textbf{Scaling mismatch}: một số nhóm feature có scale khác nhau làm lệch tối ưu
\end{itemize}

Do đó, mô hình đầy đủ không đồng nghĩa với mô hình tốt nhất,
mà cần có cơ chế chọn lọc đặc trưng và regularization phù hợp.

%------------------------------------------------
\paragraph{6. Nhận xét tổng quát}

Từ toàn bộ kết quả ablation study, có thể kết luận:

\begin{itemize}
    \item Fourier basis là thành phần quan trọng nhất trong việc giảm sai số tổng thể.
    \item RBF basis giúp cải thiện độ ổn định nhưng không mạnh bằng Fourier.
    \item Polynomial basis gây mất ổn định số học khi tăng độ phức tạp.
    \item MedInc và Spatial features là hai nhóm đặc trưng quyết định chính.
    \item Mô hình tốt không phải là mô hình đầy đủ, mà là mô hình cân bằng giữa
    biểu diễn và khả năng tổng quát hóa.
\end{itemize}

Tổng thể, kết quả nhấn mạnh vai trò của \textit{feature importance under structural perturbation},
trong đó việc loại bỏ hoặc thay đổi từng thành phần giúp làm rõ cấu trúc thông tin tiềm ẩn
trong dữ liệu.

\paragraph{Kết luận}

Tổng hợp từ cả validation curve và ablation study, có thể thấy hiệu năng của mô hình
không chỉ phụ thuộc vào độ phức tạp của từng thành phần riêng lẻ, mà còn phụ thuộc mạnh
vào sự tương tác giữa các nhóm đặc trưng và không gian hàm cơ sở.

Xét từ góc nhìn bias-variance tradeoff, các mô hình thể hiện hành vi nhất quán:

\begin{itemize}
    \item Khi độ phức tạp tăng (Polynomial degree hoặc số basis tăng), bias giảm
    nhưng variance tăng mạnh, dẫn đến hiện tượng overfitting rõ rệt.

    \item Các mô hình có cấu trúc cục bộ (RBF) hoặc theo miền tần số (Fourier)
    cho khả năng kiểm soát variance tốt hơn so với Polynomial basis toàn cục.

    \item Hiệu năng tối ưu đạt được khi mô hình nằm ở vùng cân bằng giữa biểu diễn
    (representation power) và độ ổn định số học (numerical stability).
\end{itemize}

Đặc biệt, từ ablation study, có thể thấy rằng:

\begin{itemize}
    \item Các đặc trưng kinh tế (MedInc) và vị trí địa lý (Latitude, Longitude)
    đóng vai trò quyết định trong bài toán dự đoán.

    \item Một số đặc trưng có ảnh hưởng rất nhỏ (AveOccup, Density),
    cho thấy tồn tại redundancy trong không gian dữ liệu.

    \item Việc loại bỏ Fourier basis gây suy giảm nghiêm trọng hiệu năng,
    chứng tỏ dữ liệu có thành phần cấu trúc theo miền tần số thấp.
\end{itemize}

Một điểm quan trọng là mô hình full không phải lúc nào cũng tối ưu.
Hiện tượng FULL model có MSE cao hơn nhiều cấu hình rút gọn cho thấy:

\begin{itemize}
    \item Sự xuất hiện của noise amplification khi số lượng feature quá lớn
    \item Collinearity giữa các basis function khác nhau
    \item Thiếu regularization phù hợp trong không gian đặc trưng mở rộng
\end{itemize}

%------------------------------------------------
\paragraph{Kết luận}

Trong nghiên cứu này, chúng tôi đã xây dựng và đánh giá các mô hình hồi quy phi tuyến
dựa trên ba họ hàm cơ sở: Polynomial, RBF và Fourier, đồng thời thực hiện ablation study
để phân tích mức độ quan trọng của từng thành phần.

Kết quả thực nghiệm cho thấy:

\begin{itemize}
    \item Fourier basis cho hiệu năng tốt nhất và ổn định nhất trong hầu hết các cấu hình,
    đặc biệt trong việc mô hình hóa các thành phần dao động của dữ liệu.

    \item RBF basis cung cấp sự cân bằng tốt giữa độ linh hoạt và tính ổn định,
    phù hợp với các bài toán có cấu trúc phi tuyến cục bộ.

    \item Polynomial basis dễ dẫn đến mất ổn định số học và overfitting khi tăng bậc.

    \item Các đặc trưng kinh tế và vị trí địa lý là yếu tố quan trọng nhất
    trong bài toán dự đoán giá nhà.

    \item Mô hình tốt không nhất thiết là mô hình phức tạp nhất,
    mà là mô hình có cấu trúc phù hợp với phân bố dữ liệu.
\end{itemize}

Tổng thể, nghiên cứu khẳng định vai trò quan trọng của việc lựa chọn
không gian biểu diễn (feature space design) và kiểm soát độ phức tạp mô hình,
đồng thời nhấn mạnh nguyên lý cốt lõi trong machine learning:
\textit{generalization is achieved through controlled complexity, not maximal complexity}.